% =============================================================================
% Plantilla de presentación — Análisis de Señales y Sistemas (UTN FRC)
% Compilar con:  latexmk            (usa la configuración de latexmkrc)
%
% Opciones útiles de la clase:
%   aspectratio=169 | 43     formato panorámico (por defecto) o 4:3
%   handout                  versión para imprimir/compartir: colapsa las capas (\pause, <2->)
% =============================================================================
\documentclass[aspectratio=169, 11pt]{beamer}
% \documentclass[aspectratio=169, 11pt, handout]{beamer}

\usetheme{utnfrc}
\usepackage{senales}
\usepackage{appendixnumberbeamer}   % las diapositivas del apéndice no suman al total

\usepackage[backend=biber, style=numeric, sorting=nyt]{biblatex}
\addbibresource{bibliografia.bib}

% =============================================================================
% Datos de la presentación: ÚNICO archivo a editar para portada, pie y PDF
% =============================================================================
\title{Transformadas en el análisis de sistemas LTI}
\subtitle{Convolución, Fourier, Laplace y transformada Z}
\materia{Análisis de Señales y Sistemas}
\curso{Ingeniería Electrónica \textperiodcentered\ Curso 3R1}
\author[Apellido, Apellido]{Nombre Apellido \and Nombre Apellido}
\docentes{Ing.~Nombre Apellido, Ing.~Nombre Apellido}
\institute{Universidad Tecnológica Nacional \textperiodcentered\ Facultad Regional Córdoba}
\date{Noviembre de 2026}

% Metadatos del PDF
\hypersetup{
  pdftitle    = {Transformadas en el análisis de sistemas LTI},
  pdfauthor   = {Nombre Apellido, Nombre Apellido},
  pdfsubject  = {Análisis de Señales y Sistemas — UTN FRC},
  pdfkeywords = {señales, sistemas LTI, Fourier, Laplace, transformada Z},
}


\begin{document}

\portada

\begin{frame}{Contenido}
  \tableofcontents
\end{frame}

% =============================================================================
\section{Señales}
% =============================================================================

\begin{frame}{Señales: definición y clasificación}
  \begin{definicion}
    Una \alert{señal} es una función que transporta información sobre el
    comportamiento de un fenómeno físico: $x\colon \R \to \C$ (tiempo continuo)
    o $x\colon \Zent \to \C$ (tiempo discreto).
  \end{definicion}

  \begin{columns}[T]
    \begin{column}{0.48\textwidth}
      \begin{itemize}
        \item Continua $x(t)$ o discreta $x[n]$
        \item Periódica: $x(t) = x(t + T_0)\ \forall t$
        \item Par / impar: $x(-t) = \pm x(t)$
      \end{itemize}
    \end{column}
    \begin{column}{0.48\textwidth}
      \begin{itemize}
        \item Determinística o aleatoria
        \item Causal: $x(t) = 0$ para $t < 0$
        \item De energía o de potencia
      \end{itemize}
    \end{column}
  \end{columns}
\end{frame}

\begin{frame}{Señales elementales}
  \begin{columns}[c]
    \begin{column}{0.42\textwidth}
      Escalón unitario
      \[
        \escalon{t} =
        \begin{cases}
          1, & t \geq 0\\
          0, & t < 0
        \end{cases}
      \]
      Impulso unitario
      \[
        \impulso{t} = \frac{\dd \escalon{t}}{\dd t},
        \qquad \int_{-\infty}^{\infty} \impulso{t} \dd t = 1
      \]
      \pause
      Exponencial discreta causal
      \[
        x[n] = a^{n}\,\escalon{n}, \quad \abs{a} < 1
      \]
    \end{column}
    \begin{column}{0.54\textwidth}
      \centering
      \begin{tikzpicture}
        \begin{axis}[ejes senal, width=\linewidth, height=3.4cm,
            xlabel={$t$}, ylabel={$\escalon{t}$}, xmin=-2, xmax=3, ymin=0, ymax=1.2,
            xtick={-2,...,3}, ytick={1}]
          \addplot[senal continua, domain=-2:0] {0};
          \addplot[senal continua, domain=0:3] {1};
          \draw[azulutn, very thick, densely dotted] (axis cs:0,0) -- (axis cs:0,1);
        \end{axis}
      \end{tikzpicture}

      \onslide<2->{%
      \begin{tikzpicture}
        \begin{axis}[ejes senal, width=\linewidth, height=3.4cm,
            xlabel={$n$}, ylabel={$x[n]$}, xmin=-3, xmax=10, ymin=0, ymax=1.2,
            xtick={-2,0,...,10}, ytick={1}]
          \addplot[senal discreta, domain=-3:-1, samples=3] {0};
          \addplot[senal discreta, domain=0:10, samples=11] {0.75^x};
        \end{axis}
      \end{tikzpicture}}
    \end{column}
  \end{columns}
\end{frame}

\begin{frame}{Energía y potencia}
  \begin{definicion}[energía y potencia media]
    \vspace{-0.5\baselineskip}
    \begin{align*}
      E_x &= \int_{-\infty}^{\infty} \abs{x(t)}^2 \dd t, &
      P_x &= \lim_{T \to \infty} \frac{1}{2T} \int_{-T}^{T} \abs{x(t)}^2 \dd t
    \end{align*}
  \end{definicion}

  \pause
  \begin{ejemplo}
    Para $x(t) = A\,\e^{-at}\escalon{t}$, con $a > 0$:
    \[
      E_x = \int_{0}^{\infty} A^2 \e^{-2at} \dd t
          = \resaltar<3->{\frac{A^2}{2a}} < \infty
      \quad\Longrightarrow\quad P_x = 0
    \]
    \onslide<3->{Es una \alert{señal de energía}.}
  \end{ejemplo}
\end{frame}

% =============================================================================
\section{Sistemas lineales e invariantes en el tiempo}
% =============================================================================

\begin{frame}{Propiedades de los sistemas}
  Un sistema $\mathcal{T}$ transforma una entrada en una salida: $y(t) = \mathcal{T}\{x(t)\}$.

  \begin{definicion}[linealidad]
    $\mathcal{T}\{a\,x_1(t) + b\,x_2(t)\} = a\,\mathcal{T}\{x_1(t)\} + b\,\mathcal{T}\{x_2(t)\}$
    para todo $a, b \in \C$.
  \end{definicion}

  \begin{definicion}[invariancia en el tiempo]
    Si $y(t) = \mathcal{T}\{x(t)\}$, entonces $y(t - t_0) = \mathcal{T}\{x(t - t_0)\}$ para todo $t_0$.
  \end{definicion}

  \begin{importante}
    Un sistema LTI queda completamente caracterizado por su
    \alert{respuesta al impulso} $h(t) = \mathcal{T}\{\impulso{t}\}$.
  \end{importante}
\end{frame}

\begin{frame}{Integral de convolución}
  \begin{columns}[c]
    \begin{column}{0.5\textwidth}
      \begin{teorema}[convolución]
        \[
          y(t) = x(t) \conv h(t) = \int_{-\infty}^{\infty} x(\tau)\, h(t - \tau) \dd \tau
        \]
      \end{teorema}
      \begin{itemize}
        \item<2-> Reflejar $h(\tau) \to h(-\tau)$
        \item<3-> Desplazar $t$: $h(t - \tau)$
        \item<4-> Integrar el producto (área sombreada)
      \end{itemize}
    \end{column}
    \begin{column}{0.46\textwidth}
      \centering
      \begin{tikzpicture}
        \begin{axis}[ejes senal, width=\linewidth, height=4.6cm,
            xlabel={$\tau$}, xmin=-1.5, xmax=3, ymin=0, ymax=1.3,
            xtick={0,2}, extra x ticks={1.2}, extra x tick labels={$t$},
            extra x tick style={tick label style={vision}}, ytick={1}]
          % x(τ): pulso rectangular en [0, 2]
          \addplot[senal continua, const plot] coordinates {(-1.5,0) (0,1) (2,0) (3,0)};
          % h(t-τ) con h(t) = e^{-t} u(t), para t = 1,2
          \only<3->{\addplot[energia, very thick, samples=100, domain=-1.5:1.2] {exp(-(1.2-x))};
                    \addplot[energia, very thick, domain=1.2:3] {0};}
          \only<4->{\addplot[fill=celeste, fill opacity=0.35, draw=none, domain=0:1.2, samples=60]
                      {exp(-(1.2-x))} \closedcycle;}
          \node[azulutn, anchor=south, font=\small] at (axis cs:1.6,1) {$x(\tau)$};
          \only<3->{\node[energia!85!black, anchor=south east, font=\small]
                      at (axis cs:1.15,1.02) {$h(t-\tau)$};}
        \end{axis}
      \end{tikzpicture}
    \end{column}
  \end{columns}
\end{frame}

\begin{frame}{Suma de convolución}
  Para sistemas LTI de tiempo discreto:
  \[
    y[n] = x[n] \conv h[n] = \sum_{k=-\infty}^{\infty} x[k]\, h[n - k]
  \]
  \begin{columns}[c]
    \begin{column}{0.5\textwidth}
      \begin{ejemplo}
        Con $x[n] = \escalon{n} - \escalon{n-4}$ y $h[n] = 0{,}5^{n}\escalon{n}$:
        \[
          y[n] = \sum_{k=0}^{\min(n,3)} 0{,}5^{\,n-k}, \quad n \geq 0
        \]
      \end{ejemplo}
    \end{column}
    \begin{column}{0.46\textwidth}
      \centering
      \begin{tikzpicture}
        \begin{axis}[ejes senal, width=\linewidth, height=4cm,
            xlabel={$n$}, ylabel={$y[n]$}, xmin=-2, xmax=10, ymin=0, ymax=2,
            xtick={0,2,...,10}, ytick={1, 1.875}, yticklabels={$1$, $1{,}875$}]
          \addplot[senal discreta] coordinates {
            (-2,0) (-1,0) (0,1) (1,1.5) (2,1.75) (3,1.875) (4,0.9375)
            (5,0.46875) (6,0.234375) (7,0.1171875) (8,0.05859375) (9,0.029296875) (10,0.0146484375)};
        \end{axis}
      \end{tikzpicture}
    \end{column}
  \end{columns}
\end{frame}

\begin{frame}{Interconexión de sistemas LTI}
  \centering
  \begin{tikzpicture}[node distance=1.1cm and 1.3cm]
    % Serie
    \node (x1) {$x(t)$};
    \node[bloque, right=of x1] (h1) {$h_1(t)$};
    \node[bloque, right=of h1] (h2) {$h_2(t)$};
    \node[right=of h2] (y1) {$y(t)$};
    \draw[flecha] (x1) -- (h1);
    \draw[flecha] (h1) -- (h2);
    \draw[flecha] (h2) -- (y1);
    \node[right=0.6cm of y1, anchor=west] {$h(t) = h_1(t) \conv h_2(t)$};

    % Paralelo
    \node[below=2.3cm of x1] (x2) {$x(t)$};
    \node[rama, right=0.7cm of x2] (r) {};
    \node[bloque, above right=0.25cm and 0.9cm of r] (g1) {$h_1(t)$};
    \node[bloque, below right=0.25cm and 0.9cm of r] (g2) {$h_2(t)$};
    \node[sumador, right=3.2cm of r] (s) {$+$};
    \node[right=0.9cm of s] (y2) {$y(t)$};
    \draw[thick] (x2) -- (r);
    \draw[flecha] (r) |- (g1);
    \draw[flecha] (r) |- (g2);
    \draw[flecha] (g1) -| (s);
    \draw[flecha] (g2) -| (s);
    \draw[flecha] (s) -- (y2);
    \node[right=0.6cm of y2, anchor=west] {$h(t) = h_1(t) + h_2(t)$};
  \end{tikzpicture}
\end{frame}

% =============================================================================
\section{Análisis de Fourier}
% =============================================================================

\begin{frame}{Serie de Fourier}
  Toda señal periódica $x(t)$ de período $T_0$ que cumple las condiciones de Dirichlet
  admite el desarrollo
  \begin{align*}
    x(t) &= \sum_{k=-\infty}^{\infty} c_k\, \e^{\j k \omega_0 t},
    & \omega_0 &= \frac{2\pi}{T_0} \\
    % En align* usar \onslide por celda (\pause no funciona bien: amsmath mide dos veces)
    \onslide<2->{c_k} &\onslide<2->{= \frac{1}{T_0} \int_{\langle T_0 \rangle} x(t)\, \e^{-\j k \omega_0 t} \dd t}
  \end{align*}
  \begin{teorema}<3->[Parseval para señales periódicas]
    \[
      P_x = \frac{1}{T_0} \int_{\langle T_0 \rangle} \abs{x(t)}^2 \dd t
          = \sum_{k=-\infty}^{\infty} \abs{c_k}^2
    \]
  \end{teorema}
\end{frame}

\begin{frame}{Transformada de Fourier}
  \begin{definicion}
    \vspace{-0.5\baselineskip}
    \begin{align*}
      X(\j\omega) = \TF{x(t)} &= \int_{-\infty}^{\infty} x(t)\, \e^{-\j\omega t} \dd t \\
      x(t) = \iTF{X(\j\omega)} &= \frac{1}{2\pi} \int_{-\infty}^{\infty} X(\j\omega)\, \e^{\j\omega t} \dd \omega
    \end{align*}
  \end{definicion}
  Se escribe $x(t) \fourier X(\j\omega)$.
  \begin{columns}[T]
    \begin{column}{0.48\textwidth}
      \begin{itemize}
        \item $X(\j\omega) = \abs{X(\j\omega)}\, \e^{\j\angle X(\j\omega)}$
        \item Espectro de magnitud y de fase
      \end{itemize}
    \end{column}
    \begin{column}{0.48\textwidth}
      \begin{itemize}
        \item $x(t)$ real $\Rightarrow X(-\j\omega) = \conj{X}(\j\omega)$
        \item Simetría hermítica
      \end{itemize}
    \end{column}
  \end{columns}
\end{frame}

\begin{frame}{Pares y propiedades}
  \begin{columns}[T]
    \begin{column}{0.47\textwidth}
      \centering
      \renewcommand{\arraystretch}{1.25}
      \begin{tabular}{@{}cc@{}}
        \toprule
        $x(t)$ & $X(\j\omega)$ \\
        \midrule
        $\impulso{t}$ & $1$ \\
        $\e^{-at}\escalon{t}$ & $\dfrac{1}{a + \j\omega}$ \\
        $\rect\!\left(\dfrac{t}{\tau}\right)$ & $\tau \sinc\!\left(\dfrac{\omega\tau}{2\pi}\right)$ \\
        $\e^{\j\omega_0 t}$ & $2\pi\, \impulso{\omega - \omega_0}$ \\
        $\cos(\omega_0 t)$ & $\pi\left[\impulso{\omega - \omega_0} + \impulso{\omega + \omega_0}\right]$ \\
        \bottomrule
      \end{tabular}
    \end{column}
    \begin{column}{0.5\textwidth}
      \centering\small
      \renewcommand{\arraystretch}{1.25}
      \begin{tabular}{@{}lc@{}}
        \toprule
        Propiedad & Resultado \\
        \midrule
        Desplazamiento & $x(t - t_0) \fourier \e^{-\j\omega t_0} X(\j\omega)$ \\
        Modulación & $\e^{\j\omega_0 t} x(t) \fourier X(\j(\omega - \omega_0))$ \\
        Derivación & $\dfrac{\dd x}{\dd t} \fourier \j\omega\, X(\j\omega)$ \\
        Convolución & $x \conv h \fourier X(\j\omega)\, H(\j\omega)$ \\
        Producto & $x\,y \fourier \dfrac{1}{2\pi}\, X \conv Y$ \\
        \bottomrule
      \end{tabular}
    \end{column}
  \end{columns}
  \vfill
  {\footnotesize Con $\sinc(x) = \dfrac{\sin(\pi x)}{\pi x}$. Tablas completas en \cite{oppenheim1997}.}
\end{frame}

\begin{frame}{Espectro de un pulso rectangular}
  \begin{columns}[c]
    \begin{column}{0.36\textwidth}
      \[
        \rect\!\left(\frac{t}{\tau}\right) \fourier \frac{2 \sin(\omega\tau/2)}{\omega}
      \]
      \begin{itemize}
        \item Ceros en $\omega = \dfrac{2\pi k}{\tau}$, $k \neq 0$
        \item<2-> Pulso más angosto $\Rightarrow$ \alert{espectro más ancho}
      \end{itemize}
    \end{column}
    \begin{column}{0.6\textwidth}
      \centering
      \begin{tikzpicture}
        \begin{axis}[espectro, width=\linewidth, height=5.2cm,
            ylabel={$X(\j\omega)$}, xmin=-20, xmax=20, ymin=-0.35, ymax=1.1,
            xtick={-12.566,-6.283,6.283,12.566},
            xticklabels={$-\frac{4\pi}{\tau}$,$-\frac{2\pi}{\tau}$,$\frac{2\pi}{\tau}$,$\frac{4\pi}{\tau}$},
            ytick={1}, yticklabels={$\tau$},
            legend style={font=\scriptsize, draw=none, at={(1,1)}, anchor=north east}]
          \addplot[senal continua, samples=400, domain=-20:20] {2*sin(deg(x/2))/x};
          \addlegendentry{$\tau = 1$}
          \only<2->{\addplot[energia, very thick, samples=400, domain=-20:20] {2*sin(deg(x*0.25))/x};
                    \addlegendentry{$\tau = 0{,}5$}}
        \end{axis}
      \end{tikzpicture}
    \end{column}
  \end{columns}
\end{frame}

% =============================================================================
\section{Transformadas de Laplace y Z}
% =============================================================================

\begin{frame}{Transformada de Laplace}
  \begin{definicion}[transformada bilateral]
    \[
      X(s) = \TL{x(t)} = \int_{-\infty}^{\infty} x(t)\, \e^{-st} \dd t,
      \qquad s = \sigma + \j\omega \in \ROC
    \]
  \end{definicion}
  \begin{ejemplo}
    \vspace{-0.5\baselineskip}
    \begin{align*}
      \e^{-at}\escalon{t} &\laplace \frac{1}{s + a}, & \Re\{s\} &> -a \\
      -\e^{-at}\escalon{-t} &\laplace \frac{1}{s + a}, & \Re\{s\} &< -a
    \end{align*}
  \end{ejemplo}
  \begin{importante}
    \small Misma expresión, señales distintas: la \alert{región de convergencia} es parte de la transformada.
  \end{importante}
\end{frame}

\begin{frame}{Polos, ceros y estabilidad}
  \begin{columns}[c]
    \begin{column}{0.5\textwidth}
      \[
        H(s) = \frac{s + 1}{(s + 2)(s^2 + 2s + 5)}
      \]
      \begin{itemize}
        \item Cero: $s = -1$
        \item Polos: $s = -2$, \ $s = -1 \pm \j 2$
        \item<2-> Sistema causal: $\ROC = \{\Re\{s\} > -1\}$
        \item<3-> La ROC contiene el eje $\j\omega$ $\Rightarrow$ \alert{BIBO estable}
      \end{itemize}
    \end{column}
    \begin{column}{0.46\textwidth}
      \centering
      \begin{tikzpicture}
        \begin{axis}[plano s, width=\linewidth, xmin=-3, xmax=1.5, ymin=-3, ymax=3,
            xtick={-2,-1}, ytick={-2,2}]
          \only<2->{\fill[celeste, opacity=0.25] (axis cs:-1,-3) rectangle (axis cs:1.5,3);
                    \draw[celeste, thick, dashed] (axis cs:-1,-3) -- (axis cs:-1,3);
                    \node[celeste!60!black, font=\small] at (axis cs:0.75,-2.4) {ROC};}
          \addplot[cero] coordinates {(-1,0)};
          \addplot[polo] coordinates {(-2,0) (-1,2) (-1,-2)};
        \end{axis}
      \end{tikzpicture}
    \end{column}
  \end{columns}
\end{frame}

\begin{frame}{Sistemas realimentados}
  \centering
  \begin{tikzpicture}[node distance=1cm and 1.4cm]
    \node (x) {$X(s)$};
    \node[sumador, right=of x] (s) {};
    \node[bloque, right=of s] (g) {$G(s)$};
    \node[rama, right=1.2cm of g] (r) {};
    \node[right=of r] (y) {$Y(s)$};
    \node[bloque, below=0.8cm of g] (h) {$H(s)$};
    \draw[flecha] (x) -- node[pos=0.85, above, senial] {$+$} (s);
    \draw[flecha] (s) -- node[above, senial] {$E(s)$} (g);
    \draw[thick] (g) -- (r);
    \draw[flecha] (r) -- (y);
    \draw[flecha] (r) |- (h);
    \draw[flecha] (h) -| node[pos=0.9, left, senial] {$-$} (s);
  \end{tikzpicture}

  \vspace{1em}
  \pause
  \begin{columns}[T]
    \begin{column}{0.44\textwidth}
      \begin{align*}
        E(s) &= X(s) - H(s)\,Y(s) \\
        Y(s) &= G(s)\,E(s)
      \end{align*}
    \end{column}
    \begin{column}{0.44\textwidth}
      \[
        \Rightarrow\quad \resaltar<3->{T(s) = \frac{Y(s)}{X(s)} = \frac{G(s)}{1 + G(s)\,H(s)}}
      \]
    \end{column}
  \end{columns}
\end{frame}

\begin{frame}{Respuesta en frecuencia: diagrama de Bode}
  \framesubtitle{Sistema de segundo orden
    $H(s) = \dfrac{\omega_n^2}{s^2 + 2\zeta\omega_n s + \omega_n^2}$, con $\omega_n = 1$}
  \centering
  \begin{tikzpicture}
    \begin{semilogxaxis}[bode, name=mag, width=0.48\textwidth, height=4.6cm,
        domain=0.1:10, ylabel={$\abs{H(\j\omega)}$ [\si{dB}]}, ymin=-40, ymax=15,
        legend pos=south west]
      \foreach \z in {0.1, 0.3, 0.7} {
        \addplot+[] {-10*log10((1 - x^2)^2 + (2*\z*x)^2)};
      }
      \legend{$\zeta = 0{,}1$, $\zeta = 0{,}3$, $\zeta = 0{,}7$}
    \end{semilogxaxis}
    \begin{semilogxaxis}[bode, at={(mag.east)}, anchor=west, xshift=1.6cm,
        width=0.48\textwidth, height=4.6cm,
        domain=0.1:10, ylabel={$\angle H(\j\omega)$ [\si{\degree}]},
        ymin=-180, ymax=0, ytick={-180,-135,-90,-45,0}]
      \foreach \z in {0.1, 0.3, 0.7} {
        \addplot+[] {-atan2(2*\z*x, 1 - x^2)};
      }
    \end{semilogxaxis}
  \end{tikzpicture}
\end{frame}

\begin{frame}{Transformada Z}
  \begin{columns}[c]
    \begin{column}{0.54\textwidth}
      \begin{definicion}
        \[
          X(z) = \TZ{x[n]} = \sum_{n=-\infty}^{\infty} x[n]\, z^{-n}
        \]
      \end{definicion}
      \[
        a^n \escalon{n} \ztransf \frac{1}{1 - a z^{-1}}, \quad \abs{z} > \abs{a}
      \]
      \begin{itemize}
        \item Polos en $z = 0{,}5$ y $z = 0{,}8\,\e^{\pm\j\pi/4}$
        \item Causal y estable $\Leftrightarrow$ todos los polos dentro del círculo unitario
      \end{itemize}
    \end{column}
    \begin{column}{0.42\textwidth}
      \centering
      \begin{tikzpicture}
        \begin{axis}[plano z, width=\linewidth, xmin=-1.4, xmax=1.4, ymin=-1.4, ymax=1.4,
            xtick={-1,1}, ytick={-1,1}]
          \fill[celeste, opacity=0.25, even odd rule]
            (axis cs:-1.4,-1.4) rectangle (axis cs:1.4,1.4)
            (axis cs:0,0) circle[radius=0.8];
          \addplot[cero] coordinates {(0,0)};
          \addplot[polo] coordinates {(0.5,0) (0.5657,0.5657) (0.5657,-0.5657)};
          \node[celeste!60!black, font=\small] at (axis cs:-1.05,1.15) {ROC};
        \end{axis}
      \end{tikzpicture}
    \end{column}
  \end{columns}
\end{frame}

% =============================================================================
\section{Conclusiones}
% =============================================================================

\begin{frame}{Conclusiones}
  \begin{enumerate}
    \item Un sistema LTI queda definido por su respuesta al impulso $h(t)$ o $h[n]$.
    \item<2-> Las transformadas convierten la convolución en producto:
      \[
        y = x \conv h \quad\Longleftrightarrow\quad Y = X \cdot H
      \]
    \item<3-> La ubicación de los polos respecto de la ROC determina causalidad y estabilidad.
  \end{enumerate}

  \onslide<4->{%
  \begin{observacion}
    Estos resultados se desarrollan en detalle en \cite{oppenheim1997, haykin2003, lathi2009}.
  \end{observacion}}
\end{frame}


\begin{frame}[allowframebreaks]{Bibliografía}
  \nocite{*}
  \printbibliography[heading=none]
\end{frame}

\diapositivacierre[nombre.apellido@frc.utn.edu.ar]{¿Preguntas?}

% ------------------------------------------- Diapositivas de respaldo ---
\appendix
% =============================================================================
% Diapositivas de respaldo (después de \appendix): no se numeran en el total
% =============================================================================

\begin{frame}{Respaldo: propiedad de convolución}
  \begin{teorema}
    Si $x(t) \fourier X(\j\omega)$ y $h(t) \fourier H(\j\omega)$, entonces
    $x(t) \conv h(t) \fourier X(\j\omega)\, H(\j\omega)$.
  \end{teorema}
  \begin{demostracion}
    \vspace{-\baselineskip}
    \begin{align*}
      \TF{x \conv h}
        &= \int_{-\infty}^{\infty} \left[ \int_{-\infty}^{\infty} x(\tau)\, h(t - \tau) \dd\tau \right]
           \e^{-\j\omega t} \dd t \\
        &= \int_{-\infty}^{\infty} x(\tau) \left[ \int_{-\infty}^{\infty} h(t - \tau)\,
           \e^{-\j\omega t} \dd t \right] \dd\tau
         = \int_{-\infty}^{\infty} x(\tau)\, \e^{-\j\omega\tau} H(\j\omega) \dd\tau \\
        &= X(\j\omega)\, H(\j\omega) \qedhere
    \end{align*}
  \end{demostracion}
\end{frame}


\end{document}
